\scanpage{28}
b) Contexte des familles \unsure{finies} de modules quadratiques : $=$ formule, avec $=$ \unsure{démonstration|démonstrati.}

c) \diagram{\begin{tikzcd}[ampersand replacement=\&]
\deleted{\unsure{I}}J \arrow[r,"p"] \& I
\end{tikzcd}}
application \inserted{(surjection \unsure{si on veut})} d'ens \deleted{\unsure{C|f}}finis,
\[
 I^{-}=\{i\in I\mid\operatorname{Card}p^{-1}(i)\ \text{impair}\}.
\]
\deleted{\unsure{Je dis que}} [p. ex $J$ ens avec relation d'équivalence \inserted{$\mathcal R$}, et $I=J/\mathcal R$]. \deleted{\unsure{O}}Alors
\[
 \unsure{\omega|w}_J\simeq\omega_{I^{-}}\wedge\bigwedge_{i\in I}\omega_{J_{\unsure{i|2}}}
\]
% slanted
\margin{\illegible{1}}

\uline{Compatibilités} entre les \unsure{isom|isom.|trois} \unsure{can.|cas.} a) b) c)

\figure{\begin{tikzpicture}
\draw (0,0) .. controls (.32,.03) and (.92,-.04) .. (1.2,-.01);
\end{tikzpicture}}

\uline{Cas d'extensions}
\[
 \{\unsure{0|\cdot}\}=\deleted{\unsure{0}}E_0\subset E_1\subset E_2\ \cdots\ \subset E_d=E
\]
\deleted{\unsure{On}} Définir un \unsure{isom.} canonique
\[
 \omega_E\simeq\bigwedge_{1\leq i\leq d}\omega_{E_i\unsure{/|{}}E_{i-1}}
\]
\figure{\begin{tikzpicture}
\draw (0,0) .. controls (.76,.06) and (1.95,.16) .. (2.90,.17);
\draw (1.04,-.05) .. controls (1.50,-.04) and (2.10,-.09) .. (2.46,-.05);
\end{tikzpicture}}
% archive sheet number: 27
